% ===== PRÉAMBULE CANONIQUE — à coller VERBATIM en tête de CHAQUE .tex =====
\documentclass[11pt,a4paper]{article}
\usepackage[utf8]{inputenc}\usepackage[T1]{fontenc}\usepackage{lmodern}
\usepackage[french]{babel}
\usepackage{amsmath,amssymb,mathtools}\usepackage{icomma}
\usepackage{siunitx}\sisetup{locale=FR}  % \qty{}{}, \si{}, \ang{}
\usepackage{xcolor}\usepackage{colortbl}
\usepackage{tikz}\usepackage{pgfplots}\pgfplotsset{compat=1.18}
\usepackage[siunitx,europeanresistors]{circuitikz} % circuits (élec.)
\usepackage[most]{tcolorbox}\usepackage{enumitem}\usepackage{array,booktabs}
\usepackage[a4paper,margin=2.2cm]{geometry}
\usetikzlibrary{arrows.meta,calc,angles,quotes,positioning,patterns,%
  backgrounds,shapes.geometric,decorations.pathmorphing,decorations.markings,intersections}
\definecolor{primary}{RGB}{222,49,99}\definecolor{primarydark}{RGB}{150,22,55}
\definecolor{defcol}{RGB}{0,114,178}\definecolor{methcol}{RGB}{52,92,148}
\definecolor{excol}{RGB}{0,135,90}\definecolor{attncol}{RGB}{200,40,55}
\tcbset{boxbase/.style={enhanced,breakable,boxrule=0.9pt,arc=2.5pt,
  left=3mm,right=3mm,top=2.3mm,bottom=2.3mm,fonttitle=\bfseries,coltitle=white}}
% --- boîtes TUTORIEL (noms EXACTS, ne pas renommer) ---
\newtcolorbox{cle}[1][]{boxbase,colback=defcol!6,colframe=defcol,
  title={\textsc{À retenir}\ifx\relax#1\relax\relax\else\textmd{ : \itshape #1}\fi}}
\newtcolorbox{methode}[1][]{boxbase,colback=methcol!7,colframe=methcol,sharp corners,
  title={\textsc{Méthode}\ifx\relax#1\relax\relax\else\textmd{ --- \itshape #1}\fi}}
\newtcolorbox{exemplebox}[1][]{boxbase,colback=excol!5,colframe=excol,
  title={\textsc{Exemple}\ifx\relax#1\relax\relax\else\textmd{ : \itshape #1}\fi}}
\newtcolorbox{piege}[1][]{boxbase,colback=attncol!6,colframe=attncol,
  title={\textsc{Piège}\ifx\relax#1\relax\relax\else\textmd{ : \itshape #1}\fi}}
\newtcolorbox{remarque}[1][]{boxbase,colback=black!4,colframe=black!45,
  title={\textsc{Remarque}\ifx\relax#1\relax\relax\else\textmd{ : \itshape #1}\fi}}
% --- boîtes EXERCICE / CORRIGÉ (noms EXACTS, auto-comptées) ---
\newtcolorbox[auto counter]{exo}[1][]{boxbase,colback=primary!4,colframe=primary,
  title={\textsc{Exercice}~\thetcbcounter\ifx\relax#1\relax\relax\else\textmd{ --- \itshape #1}\fi}}
\newtcolorbox[auto counter]{corrige}[1][]{boxbase,colback=excol!4,colframe=excol,
  title={\textsc{Corrigé}~\thetcbcounter\ifx\relax#1\relax\relax\else\textmd{ --- \itshape #1}\fi}}
\newcommand{\vv}[1]{\vec{#1}}\newcommand{\ds}{\displaystyle}
\newcommand{\R}{\mathbb{R}}\newcommand{\N}{\mathbb{N}}\newcommand{\Z}{\mathbb{Z}}
% =========================================================================
\begin{document}
\usetikzlibrary{babel}
\begin{corrige}[retrouver la longueur]
On isole $L$ dans $R=\rho L/S$ : $L=\dfrac{R\,S}{\rho}$.
\[ L=\frac{2\times10^{-6}}{1{,}7\times10^{-8}}\approx\qty{118}{\metre}. \]
\end{corrige}

\begin{corrige}[comparer deux fils]
Comme $R=\rho L/S$ (même $\rho$) :
\[ \frac{R_2}{R_1}=\frac{L_2}{L_1}\times\frac{S_1}{S_2}=3\times\frac{1}{2}=1{,}5. \]
Donc $R_2=1{,}5\times R_1=1{,}5\times4=\qty{6}{\ohm}$. (L'allongement augmente $R$, mais la section
plus grande l'atténue en partie.)
\end{corrige}

\begin{corrige}[la résistance double avec la température]
\begin{enumerate}[label=\alph*)]
  \item $\dfrac{R(T)}{R_0}=1+\alpha T=1+4\times10^{-3}\times250=1+1=2$.
  \item $R(250\,\si{\degreeCelsius})=2\times R_0=2\times50=\qty{100}{\ohm}$ : la résistance est
        \textbf{doublée}. (C'est pourquoi le filament d'une lampe, très chaud, est bien plus résistant
        qu'à froid.)
\end{enumerate}
\end{corrige}

\begin{corrige}[classer des matériaux]
\begin{enumerate}[label=\alph*)]
  \item Le meilleur conducteur a la \emph{plus petite} résistivité :
        \[ \text{cuivre} \;>\; \text{fer} \;>\; \text{carbone}. \]
  \item La conductivité $\sigma=1/\rho$ est maximale pour la plus petite $\rho$ : c'est le
        \textbf{cuivre}.
\end{enumerate}
\end{corrige}

\begin{corrige}[identifier un matériau]
\begin{enumerate}[label=\alph*)]
  \item $\rho=\dfrac{R\,S}{L}=\dfrac{0{,}85\times2\times10^{-6}}{100}=\qty{1.7e-8}{\ohm\metre}$.
  \item Cette valeur correspond au \textbf{cuivre} ($1{,}7\times10^{-8}\,\si{\ohm\metre}$).
\end{enumerate}
\end{corrige}

\end{document}
